\documentclass[conference]{IEEEtran}

\usepackage[utf8]{inputenc}
\usepackage[T1]{fontenc}
\usepackage{amsmath,amsfonts}
\usepackage{graphicx}
\usepackage{booktabs}
\usepackage{listings}
\usepackage{xcolor}
\usepackage{hyperref}

\lstset{
  basicstyle=\ttfamily\footnotesize,
  breaklines=true,
  frame=single,
  columns=fullflexible,
  keepspaces=true
}

\title{Deployment e Ottimizzazione di una Kolmogorov-Arnold Network su RISC-V simulato con gem5}

\author{%
\IEEEauthorblockN{Giovanni Vaccaro}
\IEEEauthorblockA{Progetto KAN su RISC-V/gem5\\
Email: inserire email}
}

\begin{document}
\maketitle

\begin{abstract}
Questo lavoro presenta il deployment di una Kolmogorov-Arnold Network (KAN) per regressione monodimensionale su una piattaforma RISC-V simulata con gem5. Il modello viene allenato esternamente in Python, esportato in formato JSON, convertito in header C e compilato come binario RISC-V statico. L'inferenza viene eseguita in gem5 in modalita' syscall emulation con CPU timing, cache L1/L2 e memoria DDR3 simulata. Il contributo principale e' l'analisi della pipeline software-hardware e l'ottimizzazione della valutazione delle B-spline: la formulazione ricorsiva di Cox-de Boor viene sostituita da una versione iterativa basata su programmazione dinamica. Infine, vengono discusse estensioni future orientate all'hardware, tra cui quantizzazione, LUT per le basi spline, scratchpad memory e istruzioni custom RISC-V per dot product packed INT8.
\end{abstract}

\section{Introduzione}
Le Kolmogorov-Arnold Networks sono reti neurali in cui la non-linearita' non e' concentrata in una funzione di attivazione fissa applicata ai nodi, ma in funzioni univariate apprese sugli archi. Questa caratteristica rende l'inferenza diversa da quella di un MLP tradizionale: il costo dominante non e' soltanto una moltiplicazione matrice-vettore, ma la valutazione ripetuta di funzioni spline associate agli edge.

L'obiettivo del progetto e' duplice. Primo, eseguire correttamente l'inferenza di una KAN allenata in Python su un binario RISC-V simulato con gem5. Secondo, analizzare e ottimizzare la computazione dal punto di vista hardware, con particolare attenzione a memoria, cache, rappresentazione dei coefficienti spline e struttura della valutazione B-spline.

\section{Background matematico}
\subsection{Struttura di una KAN}
In un layer KAN, l'output del nodo $j$ puo' essere scritto come somma di funzioni univariate applicate agli input:
\begin{equation}
    y_j = \sum_{i=1}^{n} \phi_{i,j}(x_i),
\end{equation}
dove $x_i$ e' l'input del layer e $\phi_{i,j}$ e' la funzione appresa associata all'arco tra il nodo di input $i$ e il nodo di output $j$.

Nel caso implementato, la rete approssima la funzione target
\begin{equation}
    f(x) = \sin(2\pi x) + 0.35\sin(10\pi x),
\end{equation}
su punti uniformemente distribuiti in $[0,1]$.

\subsection{Parametrizzazione tramite B-spline}
Ogni funzione edge e' rappresentata come combinazione lineare di B-spline:
\begin{equation}
    \phi_{i,j}(x) = \sum_{k=0}^{m-1} c_{i,j,k} B_{k,p}(x),
\end{equation}
dove $c_{i,j,k}$ sono i coefficienti appresi, $B_{k,p}$ e' la $k$-esima B-spline di grado $p$, e $p=3$ nel caso cubico.

La ricorrenza di Cox-de Boor e' definita da:
\begin{equation}
B_{i,0}(x) =
\begin{cases}
1, & t_i \le x < t_{i+1},\\
0, & \text{altrimenti},
\end{cases}
\end{equation}
\begin{equation}
B_{i,p}(x) =
\frac{x-t_i}{t_{i+p}-t_i}B_{i,p-1}(x)
+
\frac{t_{i+p+1}-x}{t_{i+p+1}-t_{i+1}}B_{i+1,p-1}(x).
\end{equation}

Per una spline cubica, al piu' quattro basi sono attive localmente. Se $x$ cade nell'intervallo $[t_k,t_{k+1})$, il valore della funzione puo' essere calcolato combinando solo una finestra locale di coefficienti:
\begin{equation}
    \phi(x) \approx c_0 b_0(x) + c_1 b_1(x) + c_2 b_2(x) + c_3 b_3(x).
\end{equation}
Questa struttura e' rilevante per l'hardware, perche' puo' essere trasformata in un dot product a quattro elementi.

\section{Pipeline software e simulazione}
\subsection{Esportazione del modello}
Il training non avviene nella repository. Il modello viene allenato esternamente e salvato come JSON. Il file contiene topologia, knot vector, control points, scale factors e mask. Lo script \texttt{json\_to\_header.py} genera \texttt{include/kan\_model.h}, che contiene array C statici con i parametri del modello.

I dati principali generati sono:
\begin{itemize}
    \item \texttt{KAN\_LAYER\_KNOTS}: knot vector per ogni layer e input;
    \item \texttt{KAN\_LAYER\_CONTROL\_POINTS}: coefficienti spline per ogni edge;
    \item \texttt{KAN\_LAYER\_SCALE\_SP}: scala della componente spline;
    \item \texttt{KAN\_LAYER\_MASK}: maschera degli edge attivi;
    \item dimensioni dei layer e metadati globali.
\end{itemize}

Poiche' questi array sono dichiarati \texttt{static const}, vengono incorporati nel binario RISC-V. Durante l'esecuzione non viene letto alcun JSON: il binario contiene gia' tutti i parametri.

\subsection{Compilazione ed esecuzione in gem5}
Lo script di build compila i file \texttt{main.c}, \texttt{bspline.c} e \texttt{kan\_inference.c} con cross-compiler RISC-V e linking statico:
\begin{lstlisting}[language=bash]
riscv64-linux-gnu-gcc -O2 -static -Wall -Wextra -Iinclude \
  src/main.c src/bspline.c src/kan_inference.c \
  -lm -o build/riscv/kan_demo_riscv
\end{lstlisting}

La simulazione usa gem5 in modalita' syscall emulation. La configurazione principale include una CPU timing RISC-V, una cache L1 istruzioni, una cache L1 dati, una L2 privata e una memoria DDR3 simulata.

\subsection{Pipeline di inferenza}
Il programma principale genera $N$ input uniformemente distribuiti in $[0,1]$. Per ogni input calcola:
\begin{enumerate}
    \item predizione KAN tramite \texttt{kan\_infer(x)};
    \item valore target tramite la funzione analitica;
    \item errore, errore assoluto, MSE, MAE e checksum.
\end{enumerate}

All'interno di \texttt{kan\_infer}, due buffer locali sullo stack, \texttt{current} e \texttt{next}, memorizzano le attivazioni del layer corrente e successivo. Per ogni layer, per ogni nodo di output, il codice scorre tutti gli input del layer. Se la mask dell'edge e' zero, l'edge viene saltato. Altrimenti viene chiamata \texttt{bspline\_eval} usando l'attivazione corrente, i knots dell'input corrente e i control points dell'edge.

La pipeline logica e':
\begin{lstlisting}
current[0] = x
for layer:
  for output_index:
    sum = 0
    for input_index:
      if mask == 0: continue
      spline = bspline_eval(current[input_index], knots, control_points)
      edge_value = scale_sp * spline
      sum += mask * edge_value
    next[output_index] = sum
  current = next
return current[0]
\end{lstlisting}

\subsection{Dove sono caricati i dati}
I coefficienti e i knots sono parte della sezione read-only del binario RISC-V. In gem5, quando il processo viene caricato in SE mode, codice, dati read-only, stack e heap vengono mappati nello spazio di memoria del processo simulato. Gli accessi ai parametri avvengono tramite normali load RISC-V. La gerarchia attraversata e':
\begin{equation}
\text{registro} \leftarrow \text{L1D} \leftarrow \text{L2} \leftarrow \text{DDR3 simulata}.
\end{equation}
Gli array \texttt{current} e \texttt{next} sono temporanei locali, quindi risiedono nello stack del processo e sono anch'essi serviti dalla gerarchia cache/memoria. Non e' presente una scratchpad memory esplicita nella versione corrente.

\section{Fix e ottimizzazione delle B-spline}
\subsection{Versione precedente: valutazione ricorsiva}
Nella versione precedente, \texttt{bspline\_eval} calcolava ogni base chiamando ripetutamente \texttt{bspline\_basis}. La funzione \texttt{bspline\_basis} implementava direttamente la ricorrenza di Cox-de Boor, richiamando se stessa su grado inferiore:
\begin{lstlisting}[language=C]
y += control_points[i] * bspline_basis(
    i, degree, x, knots,
    num_knots, num_control_points, x_max);
\end{lstlisting}

Questa scelta e' corretta matematicamente, ma inefficiente: molte sotto-basi vengono ricalcolate piu' volte. Per ogni control point si ricostruisce parte dello stesso albero ricorsivo, aumentando numero di chiamate, branch, overhead di stack e istruzioni eseguite.

\subsection{Versione ottimizzata: dynamic programming iterativo}
Nel fix, \texttt{bspline\_eval} costruisce prima tutte le basi di grado zero e poi sale iterativamente fino al grado richiesto. Vengono usati due array locali:
\begin{lstlisting}[language=C]
float prev[num_intervals];
float curr[num_intervals];
\end{lstlisting}
\texttt{prev} contiene le basi del grado precedente, mentre \texttt{curr} contiene quelle del grado corrente. Per ogni grado, il codice applica la ricorrenza usando i valori gia' calcolati invece di richiamare ricorsivamente la funzione.

La nuova struttura e':
\begin{lstlisting}[language=C]
for (int i = 0; i < num_intervals; ++i)
    prev[i] = (knots[i] <= x && x < knots[i+1]) ? 1.0f : 0.0f;

for (int current_degree = 1; current_degree <= degree; ++current_degree) {
    for (int i = 0; i < basis_count; ++i) {
        curr[i] = left + right;
    }
    for (int i = 0; i < basis_count; ++i)
        prev[i] = curr[i];
}

for (int i = 0; i < num_control_points; ++i)
    y += control_points[i] * prev[i];
\end{lstlisting}

\subsection{Perche' e' migliore}
Il fix migliora il codice per quattro motivi principali:
\begin{enumerate}
    \item elimina la ricorsione, riducendo overhead di chiamata e uso dello stack;
    \item evita il ricalcolo delle stesse basi intermedie;
    \item rende i loop piu' regolari e quindi piu' favorevoli a cache, compilatore e futura vettorizzazione;
    \item prepara il codice a ulteriori ottimizzazioni come LUT, calcolo locale delle sole basi attive e packed SIMD.
\end{enumerate}

Il miglioramento non cambia il modello matematico: calcola la stessa ricorrenza di Cox-de Boor, ma in ordine bottom-up invece che top-down ricorsivo.

\section{Sviluppi hardware-oriented}
\subsection{Quantizzazione}
Una futura estensione e' la quantizzazione dei coefficienti spline. Invece di rappresentare tutti i coefficienti in floating point, si puo' usare INT8 o INT16:
\begin{equation}
    q = \mathrm{round}\left(\frac{x}{s}\right), \qquad \hat{x}=s q.
\end{equation}
L'accumulo puo' restare in INT32 o INT64 per evitare overflow. La valutazione locale di una spline cubica diventa un prodotto tra coefficienti quantizzati e basi quantizzate.

\subsection{LUT per le basi spline}
Poiche' le basi dipendono dalla posizione locale di $x$ nell'intervallo, e non dall'edge specifico, si puo' precomputare una lookup table:
\begin{equation}
    \texttt{basis\_lut[frac][4]} = [b_0,b_1,b_2,b_3].
\end{equation}
L'edge cambia solo i coefficienti $c_0,\ldots,c_3$. Quindi il runtime puo' caricare le basi dalla LUT e combinare i coefficienti dell'edge.

\subsection{Scratchpad memory}
Una scratchpad software-managed puo' contenere:
\begin{itemize}
    \item basis LUT;
    \item tile di coefficienti packed;
    \item tile di input quantizzati;
    \item accumulatori di output.
\end{itemize}
A differenza della cache, la scratchpad non usa replacement automatico: il software decide esplicitamente quali blocchi caricare e quando riusarli.

\subsection{Custom instruction KAN.DOT4}
Con coefficienti e basi INT8, quattro valori entrano in un registro RV32. Si puo' definire una istruzione custom:
\begin{equation}
\texttt{KAN.DOT4}(rs1,rs2) = c_0b_0+c_1b_1+c_2b_2+c_3b_3,
\end{equation}
dove
\begin{equation}
rs1=[c_0,c_1,c_2,c_3], \qquad rs2=[b_0,b_1,b_2,b_3].
\end{equation}
Questa istruzione sfrutta direttamente la struttura locale della spline cubica.

\section{Conclusioni}
La repository realizza una pipeline completa di deployment: export JSON, generazione header, compilazione RISC-V, simulazione gem5 e confronto metriche. Il fix principale sostituisce la valutazione ricorsiva delle B-spline con una versione iterativa che preserva la matematica ma riduce overhead e ricalcoli. La prossima fase naturale e' trasformare il kernel in un workload hardware-friendly tramite quantizzazione, LUT, scratchpad e istruzioni custom packed.

\end{document}
